\section{Fazit und Ausblick}
\label{chap:conclusion-outlook}

Die Untersuchung beantwortet die Leitfrage durch explizite Grenzen: Kontext statt fester
Wurzel, Fähigkeiten statt Produktschablone, deklarierte Adapter statt Vermutung und
Transaktionen statt unkontrollierter Dateisystemänderung. Diese Eigenschaften unterstützen
einen austauschbaren Toolingbestand, ohne Produktquellen zu vereinnahmen.

Nächste Schritte sind zusätzliche Adapterverträge, echte plattformbezogene Messläufe,
eine aktualisierte Evidence-Sammlung und ein erneuter Template-Audit bei jeder Änderung
des technischen Vertrags. Diese Arbeiten sind \Planned{} und keine Aussage über einen
bereits erreichten Zustand.
